\section{Dark Photon (Hidden Photon) Absorption}
\label{sec:darkphoton}

\subsection{Physical process}

A kinetically-mixed dark photon $A'$ of mass $m_{A'}$ couples to ordinary
electromagnetism via
\begin{equation}
\mathcal{L} \supset -\frac{\kappa}{2}\,F_{\mu\nu}^{A'}F^{\mu\nu}_\gamma\,.
\end{equation}
If $m_{A'} \geq \Egap$, the dark photon can be \textbf{absorbed}, promoting
a valence electron to the conduction band and depositing its \textbf{full
rest energy} $m_{A'}$ as a single ionization event. Unlike DM--electron
scattering, the signal is \textbf{monochromatic}:
\begin{equation}
E_\mathrm{deposit} = m_{A'}\,.
\end{equation}

\subsection{In-medium self-energy and the energy-loss function}

The dark photon's absorption rate in a medium is governed by the imaginary
part of its transverse self-energy, which in the non-relativistic,
long-wavelength ($k\to0$) limit relevant for $m_{A'}$ in the eV range
reduces to the material's optical-limit complex dielectric function
$\varepsilon(\omega) = \varepsilon_1(\omega) + i\varepsilon_2(\omega)$.
Following \citet{Hochberg2017} and \citet{Knapen2021}, the absorption rate
per unit target mass is
\begin{equation}
\boxed{
R_\mathrm{abs}(m_{A'}) = \kappa^2\,\frac{\rhochi}{\rhoT\, m_{A'}}\,
\Im\!\left[-\frac{1}{\varepsilon(m_{A'})}\right] \times \mathcal{N}
\quad\left[\text{events kg}^{-1}\,\text{yr}^{-1}\right]
}
\label{eq:Rabs}
\end{equation}
where $\mathcal{N}$ converts natural units to events/kg/yr. The quantity
\begin{equation}
\ELF(\omega) \equiv \Im\!\left[-\frac{1}{\varepsilon(\omega)}\right]
= \frac{\varepsilon_2(\omega)}{\varepsilon_1^2(\omega)+\varepsilon_2^2(\omega)}
\end{equation}
is the \textbf{energy-loss function} — dimensionless, measurable via EELS or
optical ellipsometry, and the sole material input to the rate once $\rhoT$
is fixed. The rate is therefore a \textbf{boxcar in energy}: zero except in
a bin of physical width $\delta$ centered at $m_{A'}$, with height
$R_\mathrm{abs}/\delta$, which is exactly how CCDarkSens represents it in
the shared CSV format consumed by \texttt{ChargeIonization::FoldToNe}
(Section~\ref{sec:pipeline}).

The 90\% CL upper limit on $\kappa$ from an exposure $M\cdot T$ with
$N_{90\%} = 2.44$ (Poisson, zero observed background) is
\begin{equation}
\kappa^{90\%}(m_{A'}) = \sqrt{\dfrac{N_{90\%}}{R_\mathrm{abs}(\kappa=1,\,m_{A'})\times M\cdot T}}\,.
\end{equation}

\subsection{Standard silicon case}

The DAMIC-M PRL 2025 hidden-photon limit \citep{DAMICM2025} uses Si's
\textbf{measured} optical dielectric function
(\texttt{Si\_eps\_electron\_opticallimit.dat}) — real ellipsometry/EELS
data, not a model. Key features:
\begin{itemize}[leftmargin=*]
  \item Absorption onset at the direct gap $\approx 3.4$~eV (the indirect
        gap, 1.11~eV, does not contribute to $k\to0$ dipole absorption).
  \item A sharp Si plasmon peak at $\omega\approx16$--$17$~eV,
        $\ELF\approx30$, producing a characteristic sharp \emph{minimum} in
        the exclusion curve (a plasmon is a resonant enhancement of the
        absorption rate, which strengthens — not weakens — the limit at
        that mass).
  \item Reported as the \textbf{most stringent hidden-photon limit for
        $m_{A'}=2.5$--$24$~eV} in the PRL.
\end{itemize}
This is the ``textbook'' case: no bracketing, no scissors, no synthetic
model — the material's real measured $\varepsilon(\omega)$ is used
directly.

\subsection{\SrCdSb{} (HypMat) case}
\label{sec:darkphoton-hypmat}

No optical measurement of \SrCdSb{} exists. Three model variants bracket the
unknown $\varepsilon(\omega)$, from crudest to most structured.

\subsubsection{Drude model (primary bracketing model)}

The free-carrier Drude dielectric function with a sub-gap cutoff (no
absorption below the gap, since interband transitions are forbidden there):
\begin{equation}
\varepsilon_1(\omega) = \varepsilon_\infty - \frac{\omega_p^2}{\omega^2+\gamma^2}\,,
\qquad
\varepsilon_2(\omega) = \begin{cases}
\dfrac{\omega_p^2\,\gamma}{\omega(\omega^2+\gamma^2)} & \omega \geq \Egap \\[4pt]
0 & \omega < \Egap
\end{cases}
\end{equation}
Both parameters follow from known material properties with \textbf{no
additional DFT input}.

\paragraph{Plasma frequency} $\omega_p$ from the valence electron density
($Z_\mathrm{val}=16$ per formula unit — Sr~$2s$, 2$\times$Cd~$2s$,
2$\times$Sb~5~valence electrons):
\begin{equation}
n_e = \frac{\rho\,N_A}{M_W}\,Z_\mathrm{val}
= \frac{5.76\times6.022\times10^{23}}{555.96}\times 16
\approx 9.98\times10^{22}\ \mathrm{cm}^{-3}\,,
\end{equation}
\begin{equation}
\omega_p = \sqrt{\frac{n_e\,e^2}{\varepsilon_0\,m_e}} \approx 11.7\ \mathrm{eV}\,.
\end{equation}

\paragraph{Damping} $\gamma$ from the DC conductivity
$\sigma_\mathrm{DC}=300$~$\Omega^{-1}\mathrm{cm}^{-1}$ via the Drude
relation $\sigma_\mathrm{DC} = \varepsilon_0\,\omega_p^2/\gamma$:
\begin{equation}
\gamma = \frac{\varepsilon_0\,\omega_p^2}{\sigma_\mathrm{DC}} \approx 0.062\ \mathrm{eV} \ll \omega_p\,.
\end{equation}
This narrow damping ($\gamma/\omega_p\approx0.005$) gives a sharp,
well-resolved ELF peak at the screened plasmon frequency
$\omega_s = \omega_p/\sqrt{\varepsilon_\infty}$:

\begin{center}
\begin{tabular}{@{}lcc@{}}
\toprule
Bracket & $\varepsilon_\infty$ & $\omega_s$ \\
\midrule
Unscreened & 1 & $\approx 11.7$~eV \\
Si-like screened & 12 & $\approx 3.4$~eV \\
\bottomrule
\end{tabular}
\end{center}

Both brackets are run as independent scan configs
(\texttt{hypmat\_unscreened}, \texttt{hypmat\_screened}), giving an
upper/lower rate bracket that differs by $\approx(\varepsilon_\infty)^2$ in
the weak-absorption regime.

\paragraph{What the Drude model still misses:} it is a free-carrier model
and does not capture interband transitions above the gap; it has no
$q$-dependence (only relevant for the optical/dark-photon channel, which is
exactly $k\to0$); $\varepsilon_\infty$ itself remains bracketed rather than
measured.

\subsubsection{QCDark2-ELF variant (intermediate model)}

A third variant uses Si's \textbf{measured} optical dielectric function
(same file as the standard case) with $\Egap=2.1$~eV enforced at darkelf
init time — the scissors-shifted \emph{direct} gap consistent with the
DM-e scissors calculation (Section~\ref{sec:dme-hypmat}). This trades the
featureless Drude curve for one with real (if Si-specific) interband
structure and a sharp Si plasmon near 17--19.5~eV, at the cost of the
interband structure and plasmon position belonging to Si rather than
\SrCdSb{}.

\begin{center}
\begin{tabular}{@{}p{2.8cm}p{5.3cm}p{2.7cm}p{2.7cm}@{}}
\toprule
Model key & ELF source & Plasmon & Status \\
\midrule
\texttt{hypmat\_unscreened} & Drude, $\varepsilon_\infty=1$,
  $\sigma_\mathrm{DC}=300$~$\Omega^{-1}\mathrm{cm}^{-1}$ & Sharp (corrected
  Drude version) & Primary optimistic bracket \\
\texttt{hypmat\_screened} & Drude, $\varepsilon_\infty=12$,
  $\sigma_\mathrm{DC}=300$~$\Omega^{-1}\mathrm{cm}^{-1}$ & Sharp at
  $\approx3.4$~eV & Primary conservative bracket \\
\texttt{hypmat\_qcdark2} & Si measured optical + 2.1~eV scissors onset &
  Si plasmon $\approx17$--$19.5$~eV & Intermediate — awaiting AiiDA
  wavefunctions \\
\bottomrule
\end{tabular}
\end{center}

\textbf{Correct density is now used in all three variants:}
$\rhoT=5.76$~g/cm$^3$ (previously an 8.0~g/cm$^3$ placeholder — the
$8.0\to5.76$ correction lowers the rate by a factor $5.76/8.0=0.72$,
weakening the projected $\kappa$ limit by $\sqrt{8.0/5.76}\approx1.18\times$,
since $R_\mathrm{abs}\propto1/\rhoT$ linearly).

The Drude bracket ($\varepsilon_\infty=1$ vs.\ 12) remains the dominant
systematic — a factor of 12 on $\kappa$ — and can only be narrowed by a
refractive-index measurement ($\varepsilon_\infty=n^2$).

\subsection{Numerical binning requirement for the monochromatic signal}
\label{sec:binning}

The boxcar representation of Eq.~\eqref{eq:Rabs} is exact at the CSV level
(rate density $R_\mathrm{abs}/\delta$ inside a window of physical width
$\delta=\texttt{binsize\_eV}$, zero outside). Turning that CSV into the
histogram consumed by \texttt{ChargeIonization::FoldToNe} requires
\texttt{RateTable::MakeTH1D} to resample it onto \texttt{nbins} output
bins — and because that resampling is \textbf{bin-center interpolation},
the output bin width must be fine enough to actually land inside the
$\delta$-wide boxcar at essentially every mass point on the scan grid. This
is a correctness requirement specific to the dark photon channel:
DM-electron and Migdal spectra are smooth over eV-scale features, so
bin-center sampling at \texttt{nbins}$\approx200$ introduces no significant
error there.

\paragraph{Two artifacts from under-resolving the boxcar
(\texttt{nbins=200}, bin-center).}
\begin{itemize}[leftmargin=*]
  \item \textbf{Missing low-mass coverage.} With \texttt{Emin\_eV=0.34},
  \texttt{Emax\_eV}$\approx100$, and \texttt{nbins=200}, the first bin
  center sits at $\approx0.589$~eV. Every scanned mass $m_{A'}$ between the
  physical threshold (0.34~eV) and $\approx0.589$~eV falls between bin
  centers and is never sampled inside its own boxcar $\to$ zero signal is
  assigned to masses that are, physically, fully accessible. Limit coverage
  therefore starts well above the true band-gap threshold.
  \item \textbf{$\sim$10$\times$ signal overcounting at ``lucky'' bins.}
  When a bin center \emph{does} happen to fall inside the $\delta=0.05$~eV
  boxcar, \texttt{MakeTH1D} assigns that bin the boxcar's rate density
  $R_\mathrm{abs}/\delta$ as its content. The histogram's implicit physical
  interpretation is that this density applies across the \emph{entire} bin
  width ($\approx0.498$~eV at \texttt{nbins=200}), so downstream
  integration effectively counts $R_\mathrm{abs}\times(0.498/0.05)\approx
  10\times$ the true rate. Because the 90\% CL limit scales as
  $\kappa\propto(\text{rate})^{-1/2}$, a $10\times$ rate inflation makes the
  sensitivity at those $\sim$23 ``lucky'' mass points appear $\approx3\times$
  better in $\kappa$ than physically justified — while the surrounding
  masses (missed entirely) show zero.
\end{itemize}

The net effect at \texttt{nbins=200} was a \textbf{staircase artifact}:
sharp, jagged, alternating over/under-estimates of sensitivity as a
function of $m_{A'}$, rather than the smooth curve the physics predicts. An
earlier attempt fixed this by switching \texttt{RateTable::MakeTH1D} from
bin-center sampling to trapezoid integration over each bin (correctly
captures the boxcar regardless of where it falls within a bin), but this
produced a \emph{different} artifact: \texttt{ChargeIonization::FoldToNe}
still evaluates the ionization table $P(n_e\mid E)$ at the bin center,
which — at \texttt{nbins=200} (bin half-width $\approx0.25$~eV) — can be up
to 0.25~eV away from the true mass $m_{A'}$. Near an ionization threshold
($\Egap + n\,\ehpair$), that mismatch flips events between adjacent $n_e$
bins, producing exaggerated, misaligned staircase steps. Trapezoid
integration was therefore reverted, leaving \texttt{MakeTH1D} at bin-center
sampling (current state of \texttt{src/io/RateTable.cc}).

\paragraph{Correct fix: increase \texttt{nbins} to resolve the boxcar
directly.} Setting \texttt{nbins=4000} over
\texttt{Emin\_eV=0.34}--\texttt{Emax\_eV=100} gives a bin width of
$\approx0.025$~eV — matched to the boxcar half-width
($\delta/2=0.025$~eV for $\delta=0.05$~eV) — so that:
\begin{itemize}[leftmargin=*]
  \item every scanned $m_{A'}$ has a bin center within $\approx0.012$~eV of
  the true mass $\to$ $P(n_e\mid E)$ is evaluated at essentially the
  correct energy, and the staircase steps that remain are the
  \textbf{physical} ones, landing precisely at the ionization thresholds
  $\Egap + n\,\ehpair = 0.34 + n\times1.7778$~eV (Section~\ref{sec:klein}),
  not numerical artifacts;
  \item coverage extends down to the true band-gap threshold rather than
  the first surviving bin center of a coarse grid;
  \item the $\sim$10$\times$ overcounting bias vanishes because the bin
  width now approximately matches the physical boxcar width, so
  bin-center sampling and the true integrated rate agree to good
  approximation — reproducing the same physics the (reverted)
  trapezoid-integration fix targeted, without its ionization-table
  misalignment side effect.
\end{itemize}

Six dark-current-sweep HypMat-unscreened scan configs were updated to
\texttt{nbins=4000} accordingly:
\texttt{darkphoton\_scan\_hypmat\_unscreened\_ne\_dc1e\{1,2,3,4,6\}.json}
plus \texttt{darkphoton\_scan\_hypmat\_unscreened\_ne\_1kgyr\_fresh.json}
(the sixth DC point in the sweep), along with a new Si-Mermin reference
projection, \texttt{configs/darkphoton\_scan\_si\_mermin\_ne\_proj\_nbins4000.json}
— a sibling of the existing \texttt{darkphoton\_scan\_si\_mermin\_ne\_proj.json}
(still \texttt{nbins=200}). Other dark photon configs not yet migrated to
\texttt{nbins=4000} (most \texttt{srcd2sb2\_*}, \texttt{eu5in2sb6\_*}, the
\texttt{*\_dc1e2/3\_\{1,10,100\}gyr} variants, and the \texttt{old\_grid\_*}
files) still carry \texttt{nbins=200} and inherit the same
coverage/overcounting caveats until updated — any figure drawing from those
configs should be treated as provisional.

\paragraph{Physics finding enabled by the fix — DC sweep is signal-limited
above $\sim$2~eV.} With correct normalization, the six DC levels
($10^{-1}$--$10^{-6}$~e$^-$/pix/day) produce \textbf{nearly identical}
$\kappa(m_{A'})$ curves for $m_{A'} \gtrsim 2$~eV. Dark-current background
enters the $n_e\geq2$ bins only through pileup-like combinatorics
($\propto\mathrm{DC}^2$), which is negligible at every DC level considered
once $n_e\geq2$ is populated — i.e.\ sensitivity there is set entirely by
signal statistics and exposure, not background. DC only matters on the
\textbf{low-mass plateau} (0.34--1~eV), where only the $n_e=1$ bin is
populated and single-pixel dark current ($\propto\mathrm{DC}$) is the
dominant background. This is a genuine physics result of the corrected
calculation, not an artifact — and it means DC-level comparisons for this
channel should be read as differing only near threshold, not across the
full mass range.
